\documentclass[11pt]{article}
\usepackage[margin=1.15in]{geometry}
\usepackage{amsmath,amssymb,amsthm,mathtools}
\usepackage{xcolor}
\usepackage[colorlinks=true,linkcolor=blue!60!black,citecolor=blue!60!black,urlcolor=blue!60!black]{hyperref}

\newtheorem{theorem}{Theorem}[section]
\newtheorem{proposition}[theorem]{Proposition}
\newtheorem{lemma}[theorem]{Lemma}
\newtheorem{corollary}[theorem]{Corollary}
\newtheorem{conjecture}[theorem]{Conjecture}
\newtheorem{cproposition}[theorem]{Computational Proposition}
\theoremstyle{definition}
\newtheorem{definition}[theorem]{Definition}
\theoremstyle{remark}
\newtheorem{remark}[theorem]{Remark}

\newcommand{\C}{\mathbb{C}}
\newcommand{\R}{\mathbb{R}}
\newcommand{\Z}{\mathbb{Z}}
\newcommand{\w}{\omega}
\newcommand{\dg}{\dagger}
\newcommand{\ad}{\operatorname{ad}}
\newcommand{\ob}{\mathfrak o}
\newcommand{\kerad}{\Pi_{\ker\ad_{h_0}}}

\title{Interaction Obstructions, Resonant PT Breaking,\\
and Doubled Jordan Symmetry in Fourth-Order Theories}
\author{GPT-5.6.sol\thanks{OpenAI model and principal programme author.}
\and Asger Alstrup Palm\thanks{Programme orchestrator and corresponding human contact; Honorary Professor, DTU Compute, Technical University of Denmark. \texttt{asger@area9.dk}.}}
\date{Open-repository preprint, July 2026}

\begin{document}
\maketitle

\begin{abstract}
We formulate a deformation problem for the canonical positive
pseudo-Hermitian metric of the Pais--Uhlenbeck oscillator and for its
Krein/ghost-parity alternative, and follow both into interacting
fourth-order theories.  Four results organize the paper.  (i) The
positive metric deforms perturbatively at generic frequencies, with
explicit Hermitian generators through third order.  (ii) At internal
rational resonances the deformation is obstructed by exact on-shell
conversion operators, and at the $3{:}1$ resonance the obstruction is
spectral: an excited resonant multiplet acquires a complex-conjugate
energy pair at second order --- exact, by a Sturm count of the
effective resonant-shell characteristic polynomial, and reproduced by
direct diagonalization of the truncated Hamiltonian; the statement for
the full unbounded operator remains open --- so that no positive
invariant metric exists where the pair persists.  (iii) In the
perfect-square field theory the momentum continuum turns such isolated
resonances into generic scattering shells: the branch-changing process
$H+L\to L+L$ obstructs the analytic second-order deformation of the
pointed-sector positive metric on an open subset of its shell, with an
exact value at a rational kinematic point.  (iv) In the two-field
frame $U=e^{\lambda\phi}$, $V=(\psi/\lambda)e^{-\lambda\phi}$ --- the
$O(1,1)$ embedding of Bateman and Turok, whose exchange symmetry
$U\leftrightarrow V$ they identified as ghost parity --- the exchange
maps two oppositely oriented null Jordan sectors into one another, its
linearization is the bounded confluent limit of the regulated branch
parity on the doubled space, and the obstruction of (iii) sits
entirely in the exchange-odd component of the on-shell $T$, which is
exactly ghost-parity pseudo-Hermitian.  A self-contained charge-null
lemma and regulated embedding further show that one-sided charge
nullity is an enhanced boundary property, not a property of the split
theory: the mapped vacuum obeys
$S_{UU}/S_{VV}=(\delta/2g)^2=\epsilon/g$ and becomes one-sided only on
the confluent line, with the Bateman--Turok massless theory at
$\epsilon=\mu^2=0$.  Exact split-shell Krein pseudo-unitarity and
boundary one-sidedness are therefore distinct; the two completions
are inequivalent in the interacting theory.
All finite-dimensional and modewise identities are
machine-verified; derivations of the principal results are given in the
appendices.
An independent Science Forge computation subsequently extends the discrete
resonance hierarchy to order six: the $5{:}3$ and $7{:}1$ loci vanish through
orders two--five and obstruct at order six on exactly the predicted conversion
monomials.  Those two instances are imported here as externally reproduced
exact evidence; the all-coprime hierarchy remains a conjecture.
\end{abstract}

\section*{AI authorship and computational accountability}
GPT-5.6.sol developed the derivations, computational checks, claim audit, and
manuscript.  Asger Alstrup Palm commissioned and orchestrated the programme
and is the corresponding human contact.  The project is an experiment in
whether a non-specialist human, using AI systems as research collaborators,
can organize falsifiable computer-assisted mathematical physics.  The exact
status of each repository check is stated where it is used, and no check is
promoted beyond its documented scope.  The manuscript remains subject to
expert review.

\tableofcontents

\section{Introduction}\label{sec:intro}

A companion series \cite{paperI,paperII,paperIII,paperIV} analyzed the
\emph{free} fourth-order theories --- the Pais--Uhlenbeck (PU)
oscillator, the fourth-order scalar field, and linearized quadratic
gravity --- and established a three-level structure: one complex
spectral covariance, several real forms (involutions), and several
completions.  The positive pseudo-Hermitian completion of Bender and
Mannheim \cite{BM2008PRL,BM2008PRD} and the Krein/ghost-parity
completion of Bateman and Turok \cite{BT2026} are real forms of one
free theory, distinguished at Jordan degenerations.  This paper begins
the interacting classification: which of these free structures can be
continuously deformed when nonlinear processes are switched on?

The organizing mechanism, established here first for discrete
oscillator resonances and then for continuum scattering shells, is
\begin{equation*}
\text{metric deformation}
\ \longrightarrow\
\text{on-shell cohomological obstruction}
\ \longrightarrow\
\text{spectral complexification}.
\end{equation*}

\paragraph{Results.}
For the PU oscillator with the cubic vertex $V=-iy^3$ inherited from a
real $z^3$ interaction: the deformation exists at first order for all
$\w_1>\w_2>0$ and at second and third order for generic frequencies
(Sections~\ref{sec:first}--\ref{sec:second}); it is obstructed at
second order at $\w_1=3\w_2$ by the exact on-shell conversion class
$\propto a_1a_2^{\dg3}-a_1^{\dg}a_2^3$, gauge-independently
(Section~\ref{sec:threeone}); the obstruction is spectral --- a
complex-conjugate pair appears in the $E=27\w_2$ multiplet at
$O(\zeta^2)$, certified exactly (Section~\ref{sec:spectral}); a
third-order obstruction appears at $\w_1=\tfrac32\w_2$, and a
transfer-lattice selection rule organizes the hierarchy
(Section~\ref{sec:selection}), whose first prediction is confirmed:
$5{:}1$ is unobstructed through order 3 and obstructed at order 4,
exactly on the predicted monomial
(Proposition~\ref{cprop:fiveone}); independent exact order-six calculations
find the predicted first obstructions at $5{:}3$ and $7{:}1$
(Proposition~\ref{cprop:ordersix}) --- the hierarchy is thereby sampled at
five qualitatively distinct resonances through order six; the deformation
diverges toward the
Jordan boundary, with exact leading asymptotics at first order and
measured geometric growth beyond (Section~\ref{sec:jordanscaling}).

For the perfect-square field theory
$S=-\tfrac12\int(\Box\phi+\lambda(\partial\phi)^2)^2$, regulated by a
mass splitting: the cubic vertex is exactly the K\"all\'en polynomial
(Section~\ref{sec:kallen}); an even-ghost selection rule protects the
positive metric at first order even above the ghost-decay threshold,
while the sectorwise branch parity is obstructed there by the real
decay (Section~\ref{sec:firstfield}); at second order the
branch-changing scattering $H+L\to L+L$ obstructs the analytic
deformation of the pointed-sector positive metric on an open subset of
its shell, with the exact value $401\sqrt6/39424$ at a rational
kinematic point in fully specified conventions
(Section~\ref{sec:sector}); the exact two-field rewriting exhibits the
extended theory's exchange symmetry, its two null Jordan sectors, and
the interaction-generated Jordan dynamics (Section~\ref{sec:twofield});
and the regulated branch parity, singular on one confluent chain,
converges on the doubled space to the linearization of that exchange
(Section~\ref{sec:confluent}).

\paragraph{Relation to prior work.}
The two-field $O(1,1)$ embedding, the identification of the exchange
$U\leftrightarrow V$ as ghost parity and internal charge conjugation,
the quantum embedding relating the two-field and fourth-order
Hamiltonians and $S$-matrices, and a Krein-space Born rule with
positivity through weak ghost symmetry are all due to Bateman and
Turok \cite{BT2026}; none of that structure is claimed here.
Likewise standard are: complex eigenvalues of pseudo-Hermitian
operators occurring in conjugate pairs with matching Jordan data
\cite{Mostafazadeh2002nd}, the forced indefiniteness of any compatible
metric in their presence \cite{FeinbergZnojil2021}, antilinear
(CPT-type) symmetry as the origin of real-or-conjugate-pair spectra
\cite{Mannheim2016}, complex-ghost-pair unitarity of Lee--Wick type
\cite{LiuModestoCalcagni2022} (physically distinct: there the complex
ghosts are excluded from asymptotic states), and
positive invariant graph subspaces of block $J$-self-adjoint operators
\cite{AzizovIokhvidov}.  What this paper adds is the
\emph{deformation-obstruction theory of the positive pseudo-Hermitian
metric} in the interacting theories: the obstruction complex itself,
the exact $3{:}1$ and $3{:}2$ conversion classes and the
transfer-lattice selection rule, the certified spectral pair and its
broken-PT tongue, the continuum scattering-shell obstruction, the
confluent bridge from the split branch parity to the doubled Jordan
exchange, and the resulting separation statement: the Bateman--Turok
Krein completion is \emph{not} equivalent to a hidden positive
pseudo-Hermitian completion, because the latter is obstructed by
on-shell branch-changing interactions while the doubled Krein pairing
remains exactly pseudo-unitary on the verified matrix elements
(Proposition~\ref{cprop:krein}).  The stronger one-sided-charge
condition entering the Bateman--Turok Born rule is instead a boundary
property (Proposition~\ref{cprop:embedding}); its explicit Born-trace
evaluation is the remaining capstone calculation.

\paragraph{Strength of the statements.}
Throughout we distinguish (a) \emph{analytic metric obstruction} ---
nonexistence of an order-by-order deformation of the canonical metric
in the specified representation and operator class --- from (b)
\emph{absolute nonexistence of positive metrics}.  For the oscillator
at $3{:}1$ the stronger statement (b) holds wherever the certified
complex pair persists, since a positive metric would make the
Hamiltonian similar to a self-adjoint operator; the qualifications
(formal perturbation theory, unbounded cubic interaction, truncation)
are stated where they apply.  For the field theory only (a) is
established here.  ``Generic'' is likewise used in three distinct
senses, always resolvable from context: \emph{generic oscillator
frequencies} (outside the finite resonance-candidate set at the stated
order), \emph{generic field obstruction} (nonzero on a nonempty open
subset of an allowed scattering shell), and \emph{kinematically open
for every mass pair} (existence of the shell, which does not by itself
imply nonvanishing of the obstruction at every point of it).  The
status ledger is: $R_1$ Jordan scaling --- theorem; $R_2,R_3$ scaling
--- computational propositions; $R_n\sim\epsilon^{-3n/2}$ ---
conjecture; the $5{:}1$ obstruction --- computational proposition
(confirmed at order 4, Proposition~\ref{cprop:fiveone}); the $5{:}3$ and
$7{:}1$ order-six obstructions --- clean-pinned externally reproduced
computational proposition (Proposition~\ref{cprop:ordersix}); a full-field
complex spectrum --- open; the exact doubled classical involution ---
theorem (Bateman--Turok, restated); its quantum implementation ---
open; the $\kappa$-odd localization of the shell obstruction ---
computational proposition at the rational point
(Proposition~\ref{cprop:krein}); one-sidedness of the regulated vacuum
image --- computational proposition, exact iff $\epsilon=0$ in the
stated charge frame (Proposition~\ref{cprop:embedding}); the boundary
Born-trace evaluation --- open.

\paragraph{Bookkeeping.}
The deformation order is counted by a formal parameter $\zeta$:
$H(\zeta)=H_0+\zeta H_3+\zeta^2H_4$, defining $R_1,R_2,\dots$ and the
obstruction orders.  The perfect-square coupling $\lambda$ is a fixed
parameter: in the two-field frame the quadratic Jordan structure and
the vertex coefficients carry $g=\lambda^2$, while the nonlinear field
map itself depends on $\lambda$.  In the oscillator sections $\zeta$
multiplies the cubic vertex and coincides with the coupling.

\paragraph{Verification and reproducibility.}
Every finite-dimensional, modewise, and distributional identity in this
paper is machine-verified at the root of the
\texttt{area9innovation/weyl-gravity} repository accompanying the
series (SymPy 1.14.0, NumPy 2.4.4, Python 3.14; cite the repository
commit hash to fix the artifact).  Appendix~\ref{app:checks} maps
checks to theorems and classifies each check as exact-symbolic,
exact-rational (Sturm/charpoly), or high-precision numerical.
Appendices~\ref{app:recursion}--\ref{app:parity} contain readable
derivations of the principal results.  Operator-domain caveats
(unbounded interactions, formal series) are flagged where they occur.
The order-six extension is pinned separately by
\path{bridge/science_forge/certificates/SCIENCE_FORGE_PU_ORDER6_IMPORT.json}
to Tango commit
\path{ccee11d08de28d3e2681db209264887950666e87}.  It records exact agreement
between two Forge backends and an independent SymPy recomputation, while also
recording that the upstream certificate was emitted by a dirty development
toolchain.  A second run built the compiler from a clean \texttt{git archive}
of that commit and reproduced all 31 checks and all three SymPy golden
comparisons; the clean replay receipt is part of the import.

\paragraph{Conventions.}
Oscillator conventions follow \cite{paperI,paperII}: normalized
coordinates $V=(x,y,p,q)$, mode 1 $=(x,p)$ at frequency $\w_1$, mode 2
$=(y,q)$ at $\w_2$, positive normal form
$h_0=\tfrac12[\w_1\w_2p^2+(\w_1/\w_2)x^2+(\w_2/\w_1)q^2+\w_1\w_2y^2]$,
rapidity $r=\log\frac{\w_1+\w_2}{\w_1-\w_2}$,
$c=\cosh\tfrac r2=\w_1/\sigma$, $s=\sinh\tfrac r2=\w_2/\sigma$,
$\sigma=\sqrt{\w_1^2-\w_2^2}$.  Operator computations use exact
Weyl-symbol Moyal calculus (Appendix~\ref{app:recursion}).

\section{The deformation complex}\label{sec:complex}

\begin{definition}[Deformation problem and obstruction classes]
\label{def:complex}
Let $\rho_0$ be the canonical free Dyson map,
$h_0=\rho_0H_0\rho_0^{-1}=h_0^{\dg}$, and $v=\rho_0V\rho_0^{-1}$ the
transported vertex.  Seek
$\rho_\zeta=\exp[-\tfrac12\sum_n\zeta^nR_n]\rho_0$ with
$R_n=R_n^{\dg}$ such that $h_\zeta=\rho_\zeta H_\zeta\rho_\zeta^{-1}$
is Hermitian order by order.  The first two equations are
\begin{equation}\label{eq:defeq}
[h_0,R_1]=v^{\dg}-v,
\qquad
[h_0,R_2]=(v_2^{\dg}-v_2)+\tfrac12[R_1,\,v+v^{\dg}],
\end{equation}
where $v_2$ is the transported $O(\zeta^2)$ vertex when present
(derivation: Appendix~\ref{app:recursion}).  For \emph{discrete}
spectrum, $h_0=\sum_EE\,P_E$, the equation $[h_0,R]=S$ has a Hermitian
solution iff $P_ESP_E=0$ for every $E$; the obstruction class is
$\ob_+=\kerad S=\sum_EP_ESP_E$, and the solution is
$R=\sum_{E\neq E'}P_ESP_{E'}/(E-E')$ plus an arbitrary Hermitian
element of the commutant.  The parallel Krein problem deforms the ghost
parity: $[H_0,K_1]=[\kappa_0,V]$, $\{\kappa_0,K_1\}=0$, with
$\ob_K=\Pi_{\ker\ad_{H_0}}[\kappa_0,V]$.
\end{definition}

\begin{remark}[Discrete versus continuous spectrum]\label{rem:volume}
The kernel projection above is defined for point spectrum.  All
field-theoretic computations in this paper are performed at
\emph{finite volume} (discrete momenta), where $\ad_{h_0}$ has point
spectrum and Definition~\ref{def:complex} applies verbatim; continuum
statements are then statements about the finite-volume family, with
shell conditions realized by tuning masses to lattice momenta.  A
distributional (rigged-Fock or spectral-measure) formulation of the
continuum obstruction is deferred.
\end{remark}

Anti-Hermitian (unitary) generator components are free and do not
affect the Hermiticity conditions; commutant additions to $R_n$ are the
metric ambiguity.  \emph{Gauge independence} of an obstruction class
means invariance under both, at all lower orders.

\section{The cubic PU oscillator at first order}\label{sec:first}

The real interaction $z^3$ continues, under the PU rotation $z=iy$, to
$V=-iy^3$; the transported vertex is $v=-i(cy+isp)^3$, and
$v^{\dg}-v=2ic\,(c^2y^3-3s^2yp^2)$.

\begin{theorem}[First-order deformation]\label{thm:first}
For every $\w_1>\w_2>0$ the first-order equation is solved by the
Hermitian Weyl-ordered cubic
\begin{equation}\label{eq:R1}
R_1=\mathcal W\!\left[
\frac{2c^3}{\w_1\w_2}qy^2
+\frac{4c^3}{3\w_1^3\w_2}q^3
-\frac{6cs^2(2\w_1^2-\w_2^2)}{\w_1\w_2(4\w_1^2-\w_2^2)}p^2q
+\frac{12cs^2\w_1}{\w_2(4\w_1^2-\w_2^2)}pxy
-\frac{12cs^2\w_1}{\w_2^3(4\w_1^2-\w_2^2)}qx^2
\right],
\end{equation}
with no singular denominator in the physical region; the equivalent
Hermitian interaction is
$h_1=\tfrac12(v+v^{\dg})=s(3c^2y^2p-s^2p^3)$.  (Verification:
substitute \eqref{eq:R1} into $[h_0,\mathcal W[f]]=i\mathcal
W[\{h_0,f\}]$, exact for quadratic $h_0$.)
\end{theorem}

\section{Second order: generic existence}\label{sec:second}

\begin{theorem}[Generic second-order deformation]\label{thm:secondgen}
For incommensurate frequencies the second-order obstruction
$\kerad\tfrac12[R_1,v+v^{\dg}]$ vanishes; the solved $R_2$ is
Hermitian, and the assembled
$h_\zeta=e^{-X/2}\star(h_0+\zeta v)\star e^{X/2}$,
$X=\zeta R_1+\zeta^2R_2$, is Hermitian through $O(\zeta^2)$, verified
end to end in exact Moyal calculus (including the
$\tfrac18[[h_0,R_1],R_1]$ term and the $\Lambda^3$ contribution).
$R_2$ carries the resonance denominators $(\w_1-\w_2)$ and
$(\w_1-3\w_2)$; no others can vanish for $\w_1>\w_2>0$.  (Monomial
system: Appendix~\ref{app:recursion}.)
\end{theorem}

\section{The exact $3{:}1$ obstruction}\label{sec:threeone}

\begin{theorem}[Second-order obstruction at $3{:}1$]\label{thm:threeone}
At $\w_1=3\w_2$ the first-order deformation still exists, but the
second-order class is nonzero and exact:
\begin{equation}
\ob_+^{(2)}
=\frac{27\sqrt3}{320\,\w_2^4}
\left(a_1a_2^{\dg3}-a_1^{\dg}a_2^3\right),
\end{equation}
the anti-Hermitian on-shell conversion of one mode-1 quantum into three
mode-2 quanta.  The class is gauge independent under the full relevant
first-order freedom: additions to $R_1$ of Hermitian commutant elements
including the quadratic invariants $N_1$, $N_2$, $N_1N_2$ \emph{and the
resonant quartic operators} $a_1a_2^{\dg3}+a_1^{\dg}a_2^3$ and
$i(a_1a_2^{\dg3}-a_1^{\dg}a_2^3)$, and anti-Hermitian first-order
generators:
$\kerad[G,v+v^{\dg}]=\kerad[v-v^{\dg},G]=0$ for all such $G$.
(Second-order kernel additions to $R_2$ commute with $h_0$ and cannot
affect the second-order class; they enter first at third order.)
Derivation: Appendix~\ref{app:threeone}.
\end{theorem}

There is thus no analytic second-order deformation of the canonical
positive metric at $3{:}1$, within the Weyl-algebra class, continuously
connected to the free metric.

\section{Resonant PT breaking}\label{sec:spectral}

\begin{proposition}[The antilinear symmetry]\label{prop:pt}
Let $\mathcal A=\Pi\circ\Theta$, where $\Theta$ is complex conjugation
in the $(x,y)$-Schr\"odinger representation ($x,y\mapsto x,y$,
$p,q\mapsto-p,-q$, $i\mapsto-i$) and $\Pi$ the total parity
$(x,y,p,q)\mapsto-(x,y,p,q)$.  (The symbol $\Theta$ is reserved for
this conjugation; it is distinct from the Krein parity generators
$K_n$, the effective spectral matrix $K^{(2)}$, and the momentum cutoff
$K_{\max}$.)  Then $\mathcal A$ is antiunitary, $\mathcal A^2=1$, and
\begin{equation}
[\mathcal A,\,h_0]=0,\qquad [\mathcal A,\,v]=0,
\qquad\text{hence}\qquad [\mathcal A,\,h_\zeta]=0
\ \ (\zeta\in\R):
\end{equation}
$h_0$ is even and $\Theta$-real; $v$ is odd under both $\Pi$ and
$\Theta$ separately ($\Theta v\Theta=-v$, $\Pi v\Pi=-v$), hence
invariant under their composition (machine-verified).  ``PT breaking''
below means the appearance of eigenvalues not fixed by this antilinear
symmetry, i.e.\ complex-conjugate pairs.
\end{proposition}

\begin{theorem}[Certified complex pair in the $E=27\w_2$ multiplet]
\label{thm:complexpair}
At $\w_1=3\w_2$ ($\w_2=1$), the second-order effective matrix on the
ten-dimensional shell $\{|j,27-3j\rangle\}_{j=0}^9$ is real tridiagonal
with antisymmetric off-diagonal
$K_{j,j+1}=-\tfrac{27\sqrt3}{640}\sqrt{(j{+}1)n(n{-}1)(n{-}2)}$,
$n=27-3j$, and explicit rational diagonal (Appendix~\ref{app:shell}).
Because only the products $K_{j,j+1}K_{j+1,j}$ enter, its
characteristic polynomial has \emph{rational} coefficients; an exact
Sturm count shows it has exactly eight real roots, hence exactly one
complex-conjugate pair, with certified thirty-digit enclosure
\begin{equation}
\kappa_\pm=2.448807382199383966\ldots\pm
0.224586593808125108\ldots\,i .
\end{equation}
Exact diagonalization of the truncated $h_0+\zeta v$ at $\zeta=0.02$
(two cutoffs) exhibits the pair with
$\operatorname{Im}E=\pm0.2246\,\zeta^2$ to $0.5\%$.  The lowest
resonant doublet $\{|1,0\rangle,|0,3\rangle\}$ remains real.
\end{theorem}

\begin{corollary}
Wherever the complex pair persists in the interacting spectrum, no
positive-definite invariant metric exists --- analytic in $\zeta$ or
not.  The hierarchy of claims is: (1) the complex pair in the
second-order effective resonant-shell matrix is \emph{exact} (Sturm
certificate); (2) its persistence in the interacting truncations is
supported by cutoff-stable direct diagonalization; (3) an
operator-theoretic statement for the full unbounded cubic Hamiltonian
remains open.
\end{corollary}

\begin{proposition}[Broken-PT tongues of the effective multiplet]
Within the second-order effective multiplet, detuning
$\delta_r=\w_1-3\w_2$ adds $j\,\delta_r$ to the shell diagonal, and the
spectrum is complex for
$|\delta_r|<\delta_c(\zeta)=38.4151\ldots\times\zeta^2+o(\zeta^2)$.
That an order-$n$ resonance produces an $O(\zeta^n)$ tongue is the
corresponding general perturbative expectation, not a proved
proposition.
\end{proposition}

\section{Third order and the selection rules}\label{sec:selection}

\begin{theorem}[Transfer-lattice selection rule]\label{thm:selection}
Grade monomials by energy transfer, $[h_0,M]=[(\Delta n_1)\w_1
+(\Delta n_2)\w_2]M$.  Order-$n$ objects have transfers in the $n$-fold
sumset of the vertex transfer set, total degree $\le n+2$, and degree
parity $\equiv n\bmod2$ (proof: Appendix~\ref{app:lattice}).  The
order-$n$ obstruction is supported on the finite candidate sets:
$\{2{:}1\}$ at order 1, $\{3{:}1\}$ at order 2,
$\{3{:}2,2{:}1,4{:}1\}$ at order 3, with $5{:}1$ first reachable at
order 4.  The lattice gives candidates; coefficients decide.
\end{theorem}

\begin{theorem}[Third order]\label{thm:third}
At generic frequencies the third-order obstruction vanishes and $R_3$
is Hermitian with odd total degrees $\le5$.  At $\w_1=\tfrac32\w_2$ the
class is nonzero; at $(\w_1,\w_2)=(3,2)$ its exact value is
\begin{equation}
\ob_+^{(3)}
=-\frac{117\sqrt{30}}{1120}\,i\,
\big(a_1^2a_2^{\dg3}+a_1^{\dg2}a_2^3\big),
\end{equation}
and the coefficient scales exactly as $\w_2^{-13/2}$ (verified by the
ratio $2^{-13/2}$ between $(\w_1,\w_2)=(6,4)$ and $(3,2)$), so the
general form carries the dimensional factor $\w_2^{-13/2}$.  The class
is gauge independent in the sense of Theorem~\ref{thm:threeone}.  The
candidates $2{:}1$ (orders 1 and 3) and $4{:}1$ (order 3) have
vanishing coefficients.
\end{theorem}

\begin{conjecture}[Resonance hierarchy]\label{conj:hierarchy}
The coprime ratio $\w_1/\w_2=p{:}q$ first obstructs at order $p+q-2$
when the mode-2 transfer $p$ is odd.  (Verified: $3{:}1$ at order 2,
$3{:}2$ at order 3, and $5{:}1$ at order 4 ---
Proposition~\ref{cprop:fiveone}; externally reproduced exact instances:
$5{:}3$ and $7{:}1$ at order 6 --- Proposition~\ref{cprop:ordersix}.
The mechanism excluding even mode-2 transfers and the all-coprime statement
remain open.)  The
resonance set is an expanding hierarchy, not (so far) a dense set.
\end{conjecture}

\begin{cproposition}[The $5{:}1$ obstruction at fourth order]
\label{cprop:fiveone}
At $(\w_1,\w_2)=(5,1)$ the order-2 and order-3 kernel projections
vanish (the selection rule makes $5{:}1$ unreachable below order 4)
and $R_2,R_3$ exist and are Hermitian.  The fourth-order class is
nonzero, exactly
\begin{equation}
\ob_+^{(4)}
=-\frac{203125\sqrt5}{2341011456}
\big(a_1a_2^{\dg5}-a_1^{\dg}a_2^5\big)
=-\frac{13\cdot5^{13/2}}{2^{16}\,3^6\,7^2}
\big(a_1a_2^{\dg5}-a_1^{\dg}a_2^5\big),
\end{equation}
the on-shell $1\leftrightarrow5$ conversion on the predicted monomial,
gauge independent in the sense of Theorem~\ref{thm:threeone}.  The
coefficient scales exactly as $\w_2^{-9}$ (ratio $2^{-9}$ between
$(10,2)$ and $(5,1)$), extending the dimensional series
$\ob_+^{(n)}\propto\w_2^{-(5n-2)/2}$ for $n=2,3,4$; and
$R_4=O(\epsilon^{-6})$ toward the Jordan boundary, the fourth data
point of $R_n=O(\epsilon^{-3n/2})$
(Section~\ref{sec:jordanscaling}).  The order-$n$ sources are
generated programmatically from the adjoint series of
Appendix~\ref{app:recursion}, and the same code path re-derives the
order-2 and order-3 classes exactly before running order 4.
(Checks FO1--FO9, \texttt{verify\_51\_order4.py}.)
\end{cproposition}

\begin{cproposition}[The $5{:}3$ and $7{:}1$ obstructions at sixth order]
\label{cprop:ordersix}
At $(\w_1,\w_2)=(5,3)$ and $(7,1)$, the exact kernel projections at
orders two through five vanish and the corresponding $R_2,\ldots,R_5$ are
Hermitian.  The first nonzero classes occur at order six and are
\begin{align}
\ob_{+,5:3}^{(6)}
&=\frac{3863828151875\sqrt{15}}{6463101113204736}
\big(a_1^3a_2^{\dg5}-a_1^{\dg3}a_2^5\big),\\
\ob_{+,7:1}^{(6)}
&=\frac{28633766567\sqrt7}{2656254925209600000}
\big(a_1a_2^{\dg7}-a_1^{\dg}a_2^7\big).
\end{align}
They are the predicted $3\leftrightarrow5$ and $1\leftrightarrow7$
on-shell conversion kernels, with the real coefficient and antisymmetric
combination required by even-order parity.  Rescaling $(5,3)$ to $(10,6)$
multiplies the result by exactly $2^{-14}$, matching the sampled dimensional
law $\ob_+^{(n)}\propto\w_2^{-(5n-2)/2}$ at $n=6$.

The evidence consists of a 31-check exact Forge battery on two code
generators, an independent end-to-end Hermiticity assembly through order five,
and a from-scratch SymPy recomputation agreeing term by term.  The bounded
import certificate pins the source commit and artifacts.  This proposition
does not prove the hierarchy for every coprime ratio, convergence, or
nonremovability under every further admissible canonical transformation.  The
upstream dirty-toolchain marker is retained for provenance and independently
discharged by the clean pinned replay recorded with the import.
\end{cproposition}

\section{The Jordan divergence hierarchy}\label{sec:jordanscaling}

\begin{theorem}[Exact first-order Jordan asymptotics]
\label{thm:jordanexact}
With $\w_{1,2}=\w\pm\epsilon$, the explicit $R_1$ of \eqref{eq:R1} has
the exact leading Laurent behaviour
\begin{equation}
R_1=\epsilon^{-3/2}\,\mathcal W\!\left[
\frac{xyp}{2\sqrt\w}+\frac{y^2q}{4\sqrt\w}-\frac{p^2q}{4\sqrt\w}
-\frac{x^2q}{2\w^{5/2}}+\frac{q^3}{6\w^{5/2}}
\right]+O(\epsilon^{-1/2}),
\end{equation}
so $R_1=\Theta(\epsilon^{-3/2})$, while the free generator diverges
only logarithmically.
\end{theorem}

\begin{cproposition}[Higher orders]\label{cprop:scaling}
The solved $R_2$ and $R_3$ scale as $\epsilon^{-3}$ and
$\epsilon^{-9/2}$ (measured decade exponents $-3.00$ and $-4.49$ over
$\epsilon=10^{-2}\to10^{-3}$), with no small-denominator enhancement
from $\w_1-\w_2$ at second order.  These are high-precision numerical
statements at fixed orders, not proved big-$O$ estimates.
\end{cproposition}

\begin{conjecture}
$R_n=\Theta(\epsilon^{-3n/2})$, so the deformation's radius of
convergence in $\zeta$ collapses as $\zeta_c\sim\epsilon^{3/2}$ toward
the Jordan boundary.
\end{conjecture}

\section{The perfect-square vertex}\label{sec:kallen}

The interacting theory is
$S=-\tfrac12\int(\Box\phi+\lambda(\partial\phi)^2)^2$ \cite{BT2026},
with cubic vertex $V_3\propto\Box\phi(\partial\phi)^2$ and quartic
$V_4\propto\tfrac12(\partial\phi)^4$.

\begin{proposition}[K\"all\'en form]\label{prop:kallen}
On $p_1+p_2+p_3=0$ the symmetrized cubic vertex is exactly
\begin{equation}
\mathcal V_3
=p_1^2(p_2\!\cdot\!p_3)+p_2^2(p_1\!\cdot\!p_3)+p_3^2(p_1\!\cdot\!p_2)
=\tfrac12\lambda_{\mathrm K}(p_1^2,p_2^2,p_3^2).
\end{equation}
The polynomial vertex factor therefore has a \emph{second-order zero at
the massless confluent point}: $\lambda_{\mathrm K}(0,0,0)=0$ and
$\nabla\lambda_{\mathrm K}(0,0,0)=0$, while the Hessian is nonzero.  On
the regulated decay channel,
$\lambda_{\mathrm K}(m_1^2,m_2^2,m_2^2)=m_1^2(m_1-2m_2)(m_1+2m_2)$:
the vertex vanishes exactly at threshold and turns on continuously
above it.  (A full external-state selection rule for generalized Jordan
legs would additionally require the confluent limit of leg
normalizations, shell measures, and mode factors; we do not claim it
here.)
\end{proposition}

\section{First order in the regulated field theory}\label{sec:firstfield}

Regulate by $P_\delta=(\Box+m_1^2)(\Box+m_2^2)$,
$\delta=m_1^2-m_2^2>0$, at finite volume
(Remark~\ref{rem:volume}).  In the mode expansion every leg is on its
branch shell, so on the kernel the vertex coefficient is
$\tfrac12\lambda_{\mathrm K}$ of the branch masses.  The transported
field mode, derived from the free Dyson transport and verified against
the spectral Wightman function of \cite{paperIII}
(Appendix~\ref{app:field}), is
\begin{equation}
\phi'(k)=\frac{i}{\sqrt{2\delta}}
\left[\frac{a_2+a_2^{\dg}}{\sqrt{\w_2}}
+\frac{a_1-a_1^{\dg}}{\sqrt{\w_1}}\right],
\qquad
\sigma=\sqrt\delta\ \text{($k$-independent)}.
\end{equation}

\begin{theorem}[First-order protection; sectorwise parity obstruction]
\label{thm:firstfield}
(i) \emph{Even-ghost selection rule}: the transported reality factors
make $v^{\dg}-v$ vanish on every odd-ghost-count monomial
(Appendix~\ref{app:field}).  Even-ghost shells are kinematically closed
for all $m_1>m_2>0$.  Hence $\ob_+^{(1)}=0$ identically --- below and
above the decay threshold.
(ii) The sectorwise branch parity $\kappa_0=(-1)^{N_{\mathrm{ghost}}}$
has $\ob_K^{(1)}=0$ below threshold and, above it, a nonzero class on
the open $1\to22$ shell, equal there to
$\lambda_{\mathrm K}(m_1^2,m_2^2,m_2^2)$ times the leg normalizations.
\end{theorem}

\begin{cproposition}[Split-basis scaling]
At massive confluence both $R_1$ and $K_1$ diverge as $\delta^{-3/2}$;
on massless Jordan paths $m_2^2=\alpha\delta$ the positive $R_1$ scales
as $\delta^{-1/2}$ with path-independent power, while the Krein $K_1$
is $\delta^{-1/2}$ generically and $\delta^{-3/2}$ on the exceptional
path $\alpha=2/7$.  These are statements about coefficients in the
split mode basis, which itself degenerates at $\delta=0$; a
basis-independent formulation on confluent states is deferred.
\end{cproposition}

\section{The exact two-field rewriting}\label{sec:twofield}

The two-field form below is, up to normalization and sign conventions,
the $O(1,1)$ embedding of Bateman and Turok \cite{BT2026}, who
introduced the pair $(\Omega,\Upsilon)$, identified the full $O(1,1)$
symmetry including the exchange as ghost parity and internal charge
conjugation, and constructed the quantum embedding relating the two
theories.  The theorem restates their embedding in our conventions;
the content added in this section and the next two is the
chart-transition description of the exchange, the null-Jordan-sector
and interaction-generated Jordan statements, the location of the
sectorwise obstruction inside this structure, and the confluent-parity
bridge.

\begin{theorem}[Two-field form and exchange symmetry
{\cite[adapted]{BT2026}}]\label{thm:twofield}
Introducing the auxiliary field $\psi$ and the change of variables
$U=e^{\lambda\phi}$, $V=(\psi/\lambda)e^{-\lambda\phi}$ (unit pointwise
Jacobian),
\begin{equation}
\mathcal L
=-\partial_\mu U\,\partial^\mu V+\frac{\lambda^2}{2}U^2V^2
=-\tfrac12(\partial X)^2+\tfrac12(\partial Y)^2
+\frac{\lambda^2}{8}(X^2-Y^2)^2,
\qquad X,Y=\tfrac{U\pm V}{\sqrt2}.
\end{equation}
The exchange $U\leftrightarrow V$ is an exact symmetry \emph{of the
extended two-field theory}.  In the original variables it is the
transition map between the charts $\phi_A=\tfrac1\lambda\log U$ and
$\phi_B=\tfrac1\lambda\log V$,
\begin{equation}
(\phi,\psi)\ \longmapsto\
\Big(\tfrac1\lambda\log\tfrac{\psi}{\lambda}-\phi,\ \psi\Big),
\end{equation}
defined on the chart overlap ($U\neq0\neq V$), singular at the
perturbative vacuum $\psi=0$, and requiring a branch choice in the
complexified theory.  The original $(\phi,\psi)$ chart covers only
$U\neq0$ and does not contain the mirror background: the exchange is a
global symmetry of the extension, not an everywhere-defined involution
of the original field space, and it is invisible at any finite
perturbative order around the original vacuum.
\end{theorem}

Bateman and Turok warn, and we adopt the caveat, that the two theories
are inequivalent beyond perturbation theory: the fourth-order path
integral corresponds to $U>0$, while the two-field model integrates
over all $U$ and has an indefinite kinetic form \cite{BT2026}.
Statements about the mirror background $(0,1)$ are therefore
statements about the extended theory; whether both null branches are
physical sectors of one completion is open.

\begin{proposition}[Null sectors and interaction-generated Jordan
dynamics]\label{prop:sectors}
The constant zero-action stationary backgrounds form the two null
branches $UV=0$ (whether they are vacua in a given completion is a
separate question, since the kinetic form is indefinite).  The
perturbative background $(U,V)=(1,0)$ is a nonzero null configuration;
the exchange maps it to $(0,1)$.  Expanding about $(1,0)$:
$\mathcal L^{(2)}=-\partial u\partial v+\tfrac g2v^2$,
$\mathcal L^{(3)}=g\,uv^2$, $\mathcal L^{(4)}=\tfrac g2u^2v^2$
($g=\lambda^2$), with equations $\Box v=0$, $\Box u=gv$: the exact
$\Box^2$ Jordan pair, the off-diagonal supplied by the interaction
coupling; about $(0,1)$ the mirror expansion has the oppositely
oriented chain.  A chain-preserving involution commuting with a single
Jordan block is $\pm I$ (companion lemma), so exchange between
oppositely oriented sectors is the only possible realization of a
parity-like symmetry.  In this frame the interactions are pure
potentials: $H_{\mathrm{int}}=-L_{\mathrm{int}}$ exactly, with no
Ostrogradsky--Legendre corrections.
\end{proposition}

\section{The sectorwise second-order obstruction}\label{sec:sector}

Regulate the sector-$A$ theory by
$\mathcal L_0=-\partial u\partial v-\mu^2uv+\tfrac g2v^2
+\tfrac\epsilon2u^2$, with masses $m_\pm^2=\mu^2\pm\sqrt{\epsilon g}$
and branch eigenvectors $v=\rho_bu$,
$\rho_\pm=\mp\delta/(2g)$, $\delta=m_+^2-m_-^2$ (the regulator
necessarily breaks the exchange symmetry).  All computations below set
$g=1$.  The positive-frame quantization and the complete normalization
conventions --- finite volume with unit-normalized discrete modes
$[a,a^{\dg}]=1$, unit-norm external states, lattice momenta, $m_L=4$
--- are collected in Appendix~\ref{app:field}.  The second-order
on-shell source element is computed from the exact matrix-element form
\begin{equation}\label{eq:sourceME}
\langle\mathrm{out}|\,\mathcal S^{(2)}_+\,|\mathrm{in}\rangle
=\langle v_2^{\dg}-v_2\rangle
+\sum_n\frac{\langle\mathrm{out}|v_1^{\dg}|n\rangle
\langle n|v_1^{\dg}|\mathrm{in}\rangle
-\langle\mathrm{out}|v_1|n\rangle\langle n|v_1|\mathrm{in}\rangle}
{E-E_n},
\end{equation}
derived in Appendix~\ref{app:field} from
$(v^{\dg}-v)(v+v^{\dg})+(v+v^{\dg})(v^{\dg}-v)=2(v^{\dg}v^{\dg}-vv)$.

\begin{theorem}[Sectorwise obstruction on an open shell subset]
\label{thm:sector}
Consider the branch-changing scattering $H+L\to L+L$, kinematically
open at suitable momenta for every $m_H>m_L$.  At the rational
center-of-mass point $m_H=\tfrac32m_L$ (both incoming particles at
rest; outgoing momenta $\pm\tfrac34m_L$; in lattice units $m_L=4$,
$m_H=6$, momenta $\pm3$, all shell conditions exact), the source
element \eqref{eq:sourceME} has the exact value
\begin{equation}
\boxed{\ \mathcal M_{\mathrm{obs}}
=\frac{401\sqrt6}{39424\,g^2}
\ \xrightarrow{\ g=1\ }\ \frac{401\sqrt6}{39424}\ }
\end{equation}
(the $1/g^2$ dependence at fixed masses is verified exactly: each
vertex carries $g\rho^2\propto\delta^2/g$), with the quartic contact
piece exactly zero.  The computation is
truncation-complete (Lemma~\ref{lem:trunc}) and the element is analytic
in the shell data on the nonsingular part of the shell; hence
\emph{the second-order obstruction is nonzero on a nonempty open subset
of the shell containing this point}.  Independent tuned kinematic
points give $0.5233$, $0.3109$, $0.3970$ (numerical), consistent with
nonvanishing on a large region; a characterization of the zero set is
left open, and no claim is made for every mass pair.
\end{theorem}

\begin{lemma}[Truncation completeness]\label{lem:trunc}
For external momenta $\{k_a,k_b\}\to\{k_c,k_d\}$ pairwise distinct,
every intermediate Fock state contributing to \eqref{eq:sourceME} has
all momenta in the finite set generated by the externals and their
pairwise sums and differences ($s=k_a+k_b$, $t=k_a-k_c$, $u=k_a-k_d$):
in each contributing chain the surviving spectator particles must
coincide with external particles, which fixes every internal momentum;
free-momentum loops arise only in forward/disconnected kinematics,
excluded here.  (Proof: Appendix~\ref{app:field}; corroborated by exact
cutoff independence between $K_{\max}=4$ and $6$.)
\end{lemma}

\begin{remark}[Oscillator versus field theory]
\begin{center}
\begin{tabular}{c|c}
PU oscillator & field theory\\
\hline
discrete energy spectrum & continuous momentum spectrum\\
isolated rational resonance & generic scattering shell\\
$3{:}1$ first appears at order 2 & $H+L\to L+L$ appears at order 2\\
finite resonant multiplet & continuum of degenerate in/out states
\end{tabular}
\end{center}
The continuum turns the discrete resonance obstruction into a
scattering-shell obstruction.  The exact exchange symmetry does not
cure it: it maps the obstructed sector-$A$ construction to the mirror
sector-$B$ construction.
\end{remark}

\begin{cproposition}[Krein consistency of the obstructed element]
\label{cprop:krein}
Let $G=(-1)^{N_{\mathrm{ghost}}}$ be ghost-branch number parity in the
positive-frame Fock space, and let $\iota$ be the mode identification
$(b,k)_B\simeq(b,k)_A$ of the mirror sector.  At the rational
kinematic point of Theorem~\ref{thm:sector}, exactly:
(i) $H_B=\iota H_A\iota^{-1}$ (the mirror sector is the same theory in
the mirrored frame) and $H_A^{\dg}=GH_AG$ term by term, so the
mirror-adjoint relation $H_B=WH_A^{\dg}W^{\dg}$ holds with
$W=\iota\circ G$ --- i.e.\ the ``doubled pairing''
$\eta_{\mathrm{dbl}}=\operatorname{offdiag}(W^{\dg},W)$ is precisely
the two-sector unfolding of the ghost-parity Krein metric of
\cite{BT2026}, not additional structure.
(ii) On the degenerate shell $\{|H(0)L(0)\rangle,|L(3)L(-3)\rangle\}$
(exhaustively the whole reachable shell) the second-order on-shell $T$
satisfies $GTG=T^{\dg}$ exactly, while Hilbert Hermiticity fails: the
parity-even (diagonal) block is real, and the entire obstruction is
the parity-odd block,
$T(\mathrm{out},\mathrm{in})=-401\sqrt6/78848=-\mathcal
M_{\mathrm{obs}}/2$.
(iii) Consequently the finite-time interaction-picture $S$ obeys
$S_B^{\dg}WS_A=W$ identically --- Krein pseudo-unitarity is exact on
the same matrix elements on which every positive pointed metric is
obstructed.  A positive $H$-invariant graph subspace of the doubled
space exists if and only if a positive pointed metric exists
(standard $J$-self-adjoint graph argument \cite{AzizovIokhvidov},
instantiated and machine-verified), so doubling preserves
pseudo-unitarity but cannot evade the obstruction.
(Checks DQ1--DQ9, \texttt{verify\_doubled\_theory.py}.)
\end{cproposition}

The obstruction is therefore carried exactly by the $\kappa$-odd
component of the on-shell $T$ --- the component that the
weak-ghost-symmetry positivity mechanism of \cite{BT2026} does not
use.  Whether it maps into their null orthogonal ($C$) component under
their asymptotic embedding is the natural next question
(Section~\ref{sec:discussion}).

\section{The confluent parity theorem}\label{sec:confluent}

\begin{theorem}[Confluent limit of the regulated parity]
\label{thm:confluent}
Let $P_\delta b_\pm=\pm b_\pm$ be the regulated branch parity and
$c=(b_++b_-)/2$, $d=(b_+-b_-)/\delta$ the confluent Jordan basis, so
$P_\delta c=\tfrac\delta2 d$ and $P_\delta d=\tfrac2\delta c$; in the
standard column convention on the ordered basis $(c,d)$,
\begin{equation}
P_\delta=
\begin{pmatrix}0&2/\delta\\ \delta/2&0\end{pmatrix},
\qquad P_\delta^2=1,
\end{equation}
with no bounded limit as $\delta\to0$: the regulated parity cannot
converge as a chain-preserving parity inside one Jordan sector.  On the
doubled space with the \emph{oppositely oriented} confluent
identification in the mirror sector ($c_B=(b_+-b_-)/2$,
$d_B=(b_++b_-)/\delta$), the cross parity
$b_\pm^A\leftrightarrow\pm b_\pm^B$ equals, exactly and
$\delta$-independently, the sector exchange
$(c_A,d_A)\leftrightarrow(c_B,d_B)$.  The regulated branch parity thus
converges, after doubling by the oppositely oriented Jordan sector, to
\emph{the linearization of the exact sector-exchange involution around
the two mirror backgrounds}.  (Identifying the full nonlinear quantum
operator requires a separate construction, not supplied here.)
\end{theorem}

\section{Discussion}\label{sec:discussion}

\paragraph{The mechanism.}
Away from resonances the positive pseudo-Hermitian metric deforms
locally, with explicit generators.  The failures are structural:
on-shell branch-changing processes produce gauge-independent
cohomological obstructions --- isolated rational loci for the
oscillator, generic scattering shells in the field theory; at the
oscillator resonances the obstruction is certified spectral; and
independent of resonances, the deformation loses uniformity toward
Jordan degeneration.  The sectorwise branch parity is broken by real
ghost decay above threshold; what survives at the massless
perfect-square boundary is the exact sector-exchanging involution of
the extended two-field theory, whose linearization is the bounded
confluent limit of the regulated parity on the doubled space.

\paragraph{Symmetry enhancement at the boundary.}
The massless perfect-square point is not the split theory with
parameters set to zero: the K\"all\'en vertex degenerates there, the
doubled $O(1,1)$ structure emerges, and the Jordan dynamics is
generated by the interaction about its own null backgrounds --- a
boundary-emergent structure mirroring the Weyl enhancement at the
conformal boundary of the gravity companion \cite{paperIV}.

\paragraph{Separation of the two completions.}
Proposition~\ref{cprop:krein} sharpens the free-theory picture of the
companion series: in the interacting theory the two real forms are not
merely distinguished at Jordan degenerations --- they \emph{separate}.
The doubled Krein pairing is exactly preserved at the split rational
shell ($GTG=T^{\dg}$ on shell, finite-time pseudo-unitarity), while
every analytic positive pseudo-Hermitian completion of the pointed
sector is obstructed on an open shell subset.  This already excludes
identifying that Krein pairing with an ordinary positive
pseudo-Hermitian completion.  It does \emph{not} establish the stronger
Bateman--Turok weak-ghost-symmetry positivity mechanism at split mass:
Proposition~\ref{cprop:embedding} shows that the regulated vacuum image
contains both charge signs there.  The sharpest remaining form of the
reconciliation is therefore a boundary transport question: under their
asymptotic embedding, physical operators decompose into a
ghost-symmetric part and a null orthogonal
part, and only the former contributes to their generalized Born rule
\cite{BT2026}; since the obstruction is $\kappa$-odd
(Proposition~\ref{cprop:krein}), the conjecture is that it is carried
into the null component --- obstructing every positive metric while
remaining invisible to the Krein probabilities.  Verifying this on the
computed $H+L\to L+L$ element is the concrete next calculation; the
lemma and the proposition below supply its verified ingredients.

\paragraph{The charge-null mechanism.}
Throughout this paragraph $\dg$ denotes the Krein adjoint and $\tau$
the Krein trace.  The following finite-particle lemma makes the
null-relocation mechanism self-contained (it is the statement used in
\cite{BT2026}, whose general proof is deferred to their companion
paper; for boost weights it is a three-line grading argument).

\begin{lemma}[Charge null]\label{lem:chargenull}
Let a Krein space carry an $SO^+(1,1)$ charge grading whose pairing
couples charge $q$ only to charge $-q$, and suppose the Krein adjoint
preserves charge --- as it does for boost weights, where $c$ and
$c^{\dg}$ of one field carry the \emph{same} charge.  Then
$\tau(A_q)=0$ for every $q\neq0$.  Consequently an operator
$C=\sum_{q<0}C_q$ containing only strictly negative charges satisfies
$\tau(C^{\dg}C)=0$ and $\tau(B^{\dg}C)=0$ for every neutral $B$: $C$
is null and orthogonal to the neutral sector, and the generalized Born
trace of a weakly ghost-symmetric process depends on its neutral part
alone.
\end{lemma}

\begin{proof}
Boost invariance of the trace gives
$\tau(A_q)=\tau(\alpha_\sigma A_q)=e^{q\sigma}\tau(A_q)$ for all
$\sigma$, hence $\tau(A_q)=0$ for $q\neq0$; equivalently, on a
charge-graded Fock space a grading-raising operator is traceless.
Charge preservation under the adjoint gives
$q(C_{q}^{\dg}C_{r})=q+r<0$ and $q(B^{\dg}C_{r})=r<0$ term by term.
\end{proof}

\begin{cproposition}[Regulated embedding and its charge structure]
\label{cprop:embedding}
At the rational kinematic point, per momentum sector, the regulated
split theory admits a canonical Bogoliubov map onto the cross-paired
charge basis of a neutral reference theory
$-\partial u\,\partial v-\tilde\mu^2uv$ (operators $c_U$, $c_V$ with
the only nonvanishing commutators $[c_U,c_V^{\dg}]=[c_V,c_U^{\dg}]$;
charges $q(c_U^{(\dg)})=+1$, $q(c_V^{(\dg)})=-1$, the analogue of
Bateman--Turok's $b_\Omega,b_\Upsilon$).  The pointed vacuum maps to a
Gaussian state whose squeezing components obey, exactly and for every
reference dispersion,
\begin{equation}
\frac{S_{UU}}{S_{VV}}=\Big(\frac{\delta}{2g}\Big)^{2}
=\frac{\epsilon}{g},
\qquad \delta^2=4\epsilon g,
\end{equation}
i.e.\ the charge $+2$ component is sourced by the regulator
$\tfrac\epsilon2u^2$ and the charge $-2$ component by the
interaction-generated $\tfrac g2v^2$.  The embedding is strictly
one-sided --- no positive charges in the vacuum image, the hypothesis
of Lemma~\ref{lem:chargenull} for the mapped theory --- if and only if
$\epsilon=0$: the $O(1,1)$-symmetric confluent line.  The massless
Bateman--Turok theory is the point $\epsilon=\mu^2=0$ on that line.
At confluence the matched-reference limit is
$S_{VV}\to-g/(4w^2)$, or $-1/(4w^2)$ in the $g=1$ normalization used
in the computation; at $\mu^2=0$, where $w=|\mathbf k|$, this is
precisely the coefficient structure of their Eqs.~(C5)--(C6).  At
split masses the exact ratio above rules out removing the $+2$
component by changing the reference dispersion.  Under the residual
charge-preserving Bogoliubov freedom, numerical solution of the
adaptation equations runs to a degenerate (confluent-type) frame rather
than finding a bounded adapted frame.  (Checks ON1--ON4,
\texttt{verify\_obstruction\_null.py}.)
\end{cproposition}

The proposition is the charge-frame image of the first-order result of
Section~\ref{sec:firstfield}: real ghost decay above threshold breaks
the sectorwise branch parity, and in the charge frame the same
$\epsilon$ sources the charge $+2$ contamination that violates weak
ghost symmetry.  The relocation of the $\kappa$-odd obstruction into a
null negative-charge component is therefore exact on the one-sided
confluent line, with the Bateman--Turok identification at its massless
point, and holds at split masses with charge contamination of relative
size $\epsilon/g$ (not assumed small away from the boundary).  The
remaining, precisely-posed step --- the obstruction-to-null theorem ---
is the boundary Born-trace evaluation: the neutral component's
blindness to the obstruction coefficient and its limit along a
specified on-shell family with $(\epsilon,\mu^2)\to(0,0)$.
One exact rational family is
$m_L=4s$, $m_H=6s$, $|\mathbf k_{\mathrm{out}}|=3s$,
$\mu^2=26s^2$, $\epsilon g=100s^4$, followed by $s\to0$; the process
operator must be transported along this family rather than held fixed
while only the embedding is varied.
Lemma~\ref{lem:chargenull} makes that step self-contained.

\paragraph{Open questions.}
Whether the two pointed sectors admit parity vacua
$(|0_A\rangle\pm|0_B\rangle)/\sqrt2$ or are superselected requires
finite-volume overlap and zero-mode analysis --- in the hyperbolic
polar coordinates $r=UV$, $\chi=\tfrac12\log(U/V)$ (a coordinate
presentation of the $O(1,1)$ model: boost $=$ $\chi$-shift, exchange
$=$ $\chi\mapsto-\chi$, conserved current $j=V\partial U-U\partial V
=2r\,\partial\chi$) this becomes a boundary/self-adjoint-extension
problem at the null locus $r=0$, where the conserved charge produces a
singular $q^2/r$ barrier.  The quantum implementation of the exchange
on a completion, and any probability interpretation beyond
\cite{BT2026}, remain open; so does a normal-ordered operator Ward
identity for the internal charge (the classical identity, with its
exact regulator breaking $\epsilon u(1+u)-\mu^2v$, is verified).  The
hierarchy conjecture's $5{:}1$ prediction is confirmed at fourth order
(Proposition~\ref{cprop:fiveone}), and its $5{:}3$ and $7{:}1$ instances are
confirmed at sixth order (Proposition~\ref{cprop:ordersix}); higher coprime
loci and the mechanism excluding even mode-2 transfers remain open.
These scalar questions are deferred rather than placed in the active
queue.  After the boundary Born-trace calculation, the next target is
the gravitational extension --- the cubic Einstein--Weyl vertices
against the covariant real forms of \cite{paperIV}.  The natural first
test is the physical-helicity decay $M\to h+h$: apply the uniform
massive-spin-$2$ quarter-turn and project
$v_3^{\dg}-v_3$ onto its real decay shell, with the Ward/BRST identities
checked before drawing a completion-level conclusion.  No CPT claim is made: the
exchange is a linear sector map, not a constructed antiunitary CPT
operator.

\begin{thebibliography}{12}

\bibitem{BM2008PRL}
C.~M.~Bender and P.~D.~Mannheim,
\emph{No-ghost theorem for the fourth-order derivative Pais--Uhlenbeck
oscillator model},
Phys.\ Rev.\ Lett.\ \textbf{100} (2008) 110402.

\bibitem{BM2008PRD}
C.~M.~Bender and P.~D.~Mannheim,
\emph{Exactly solvable PT-symmetric Hamiltonian having no Hermitian
counterpart},
Phys.\ Rev.\ D \textbf{78} (2008) 025022, arXiv:0804.4190.

\bibitem{BT2026}
S.~Bateman and N.~Turok,
\emph{Escape from Ostrogradsky via hidden ghost parity},
arXiv:2607.00096.

\bibitem{Mostafazadeh2005}
A.~Mostafazadeh,
\emph{Metric operator in pseudo-Hermitian quantum mechanics and the
imaginary cubic potential},
J.\ Phys.\ A \textbf{39} (2006) 10171, arXiv:quant-ph/0508195.

\bibitem{Mostafazadeh2002nd}
A.~Mostafazadeh,
\emph{Pseudo-Hermiticity for a class of nondiagonalizable
Hamiltonians},
J.\ Math.\ Phys.\ \textbf{43} (2002) 6343, arXiv:math-ph/0207009.

\bibitem{FeinbergZnojil2021}
J.~Feinberg and M.~Znojil,
\emph{Which metrics are consistent with a given pseudo-Hermitian
matrix?},
arXiv:2111.04216.

\bibitem{Mannheim2016}
P.~D.~Mannheim,
\emph{CPT symmetry without Hermiticity},
arXiv:1611.02100.

\bibitem{LiuModestoCalcagni2022}
J.~Liu, L.~Modesto and G.~Calcagni,
\emph{Quantum field theory with ghost pairs},
arXiv:2208.13536.

\bibitem{AzizovIokhvidov}
T.~Ya.~Azizov and I.~S.~Iokhvidov,
\emph{Linear Operators in Spaces with an Indefinite Metric},
Wiley, 1989.

\bibitem{paperI}
GPT-5.6.sol and A.~A.~Palm,
\emph{Canonical Positive Symplectic Diagonalization of the
Pais--Uhlenbeck Oscillator},
Paper~I of the pure-Weyl gravity programme, 2026;
file \path{paper/01-symplectic-diagonalization.pdf} in the repository
\url{https://github.com/area9innovation/weyl-gravity}.

\bibitem{paperII}
GPT-5.6.sol and A.~A.~Palm,
\emph{The Pais--Uhlenbeck Metric as a Minimum-Distortion Principle,
and the Representation Problem for the Fourth-Order Field},
Paper~II of the pure-Weyl gravity programme, 2026;
file \path{paper/02-variational-fock.pdf} in the repository
\url{https://github.com/area9innovation/weyl-gravity}.

\bibitem{paperIII}
GPT-5.6.sol and A.~A.~Palm,
\emph{The Universal Vacuum of the Fourth-Order Scalar Field:
Metric Orbits, Fock Sectors, and the Krein Boundary},
Paper~III of the pure-Weyl gravity programme, 2026;
file \path{paper/03-fourth-order-vacuum.pdf} in the repository
\url{https://github.com/area9innovation/weyl-gravity}.

\bibitem{paperIV}
GPT-5.6.sol and A.~A.~Palm,
\emph{Gauge Reduction and the Completion Problem in Fourth-Order Gravity:
PU Pairing, Covariant Real Forms, and the Conformal Jordan Boundary},
Paper~IV of the pure-Weyl gravity programme, 2026;
file \path{paper/04-fourth-order-gravity.pdf} in the repository
\url{https://github.com/area9innovation/weyl-gravity}.

\bibitem{Kato}
T.~Kato,
\emph{Perturbation Theory for Linear Operators},
Springer, 1966.

\bibitem{Heiss2012}
W.~D.~Heiss,
\emph{The physics of exceptional points},
J.\ Phys.\ A \textbf{45} (2012) 444016.

\end{thebibliography}

\appendix

\section{The Weyl--Moyal recursion}\label{app:recursion}

Operators are represented by Weyl symbols on the two canonical pairs
$(x,p)$, $(y,q)$; the star product is
$f\star g=\sum_n\frac1{n!}(i/2)^n\Lambda^n(f,g)$ with
$\Lambda=\overleftarrow\partial_x\overrightarrow\partial_p
-\overleftarrow\partial_p\overrightarrow\partial_x
+\overleftarrow\partial_y\overrightarrow\partial_q
-\overleftarrow\partial_q\overrightarrow\partial_y$, a finite sum on
polynomials.  For quadratic $h_0$,
$[h_0,\mathcal W[f]]=i\mathcal W[\{h_0,f\}]$ exactly; for cubic--cubic
commutators the $\Lambda^3$ (constant) term contributes.  Expanding
$h_\zeta=e^{-X/2}\star(h_0+\zeta v)\star e^{X/2}$,
$X=\zeta R_1+\zeta^2R_2+\zeta^3R_3$, with
$e^{-X/2}Ae^{X/2}=A+\tfrac12[A,X]+\tfrac18[[A,X],X]
+\tfrac1{48}[[[A,X],X],X]+\cdots$, the order-$\zeta^2$ and
$\zeta^3$ Hermiticity conditions are
\begin{align}
[h_0,R_2]&=(v_2^{\dg}-v_2)+\tfrac12[R_1,v+v^{\dg}],\\
[h_0,R_3]&=-(T_3-T_3^{\dg}),\\
T_3&=\tfrac18\big([[h_0,R_1],R_2]+[[h_0,R_2],R_1]\big)
+\tfrac1{48}[[[h_0,R_1],R_1],R_1]
+\tfrac12[v,R_2]+\tfrac18[[v,R_1],R_1],\notag
\end{align}
using $[h_0,R_1]=v^{\dg}-v$ and the fact that
$[v-v^{\dg},R_1]$ is Hermitian.  Equations $[h_0,R]=S$ are solved
monomial by monomial in the mode variables
$a_1^{m}a_1^{\dg n}a_2^{r}a_2^{\dg s}$, where
$[h_0,M]=\Omega_MM$ with
$\Omega_M=(n-m)\w_1+(s-r)\w_2$: $R_M=S_M/\Omega_M$ off the kernel;
the kernel component is the obstruction.

\section{The $3{:}1$ projection and gauge independence}
\label{app:threeone}

At $\w_1=3\w_2$ the second-order source
$S_2=\tfrac12[R_1,v+v^{\dg}]$ is a quartic polynomial; its kernel
monomials must satisfy $(n-m)3\w_2+(s-r)\w_2=0$ with degree $\le4$,
i.e.\ $(m-n,r-s)=\pm(1,-3)$, realized only by
$a_1a_2^{\dg3}$ and $a_1^{\dg}a_2^3$.  Direct evaluation of the two
coefficients gives
$\ob_+^{(2)}=\tfrac{27\sqrt3}{320\w_2^4}(a_1a_2^{\dg3}-a_1^{\dg}a_2^3)$.
Gauge independence: a change $R_1\to R_1+G$ with $[G,h_0]=0$ shifts
$S_2$ by $\tfrac12[G,v+v^{\dg}]$; transfer bookkeeping shows that for
every Hermitian kernel generator ($N_1,N_2$ polynomials and the
resonant quartics $a_1a_2^{\dg3}+a_1^{\dg}a_2^3$,
$i(a_1a_2^{\dg3}-a_1^{\dg}a_2^3)$) the shifted source has no kernel
component: the transfers of $v+v^{\dg}$, shifted by $\pm(1,-3)$ or
$(0,0)$, never land on $\pm(1,-3)$; machine verification confirms
$\kerad[G,v+v^{\dg}]=0$ term by term.  Anti-Hermitian generators $A_1$
enter through $\tfrac12[v-v^{\dg},A_1]$, and the same bookkeeping gives
$\kerad[v-v^{\dg},A_1]=0$.  Kernel additions to $R_2$ drop out of the
second-order equation since $[h_0,R_2^{\ker}]=0$.

\section{The effective shell matrix}\label{app:shell}

Standard second-order degenerate perturbation theory on the shell
$\mathcal E_{27}=\operatorname{span}\{|j,27-3j\rangle\}$ gives
$K^{(2)}_{ab}=\sum_{m\notin\mathcal E}v_{am}v_{mb}/(27-E_m)$ (the
first-order block vanishes: cubic $v$ cannot realize the transfer
$\pm(1,-3)$).  Evaluating the sums yields the rational diagonal
$D(n_1,n_2)=(633n_1^2-1836n_1n_2-285n_1+3969n_2^2+3051n_2+818)/26880$
and the antisymmetric off-diagonal quoted in
Theorem~\ref{thm:complexpair}.  For a tridiagonal matrix only the
products $K_{j,j+1}K_{j+1,j}\in\mathbb Q$ enter the characteristic
polynomial, computed exactly by the three-term recurrence
$p_j=(D_j-\kappa)p_{j-1}-K_{j-1,j}K_{j,j-1}p_{j-2}$; the Sturm count of
the resulting degree-10 rational polynomial is exactly 8.

\section{The transfer lattice}\label{app:lattice}

Transfer additivity: the grading by $\ad_{h_0}$ satisfies
$[h_0,AB]=[h_0,A]B+A[h_0,B]$, so transfers add under operator products,
hence under star products and commutators.  Degree bound: each
commutator of polynomial factors lowers total degree by at least two
($\Lambda^1$), so order-$n$ objects, built from $n$ cubic factors and
$n-1$ commutators (with $[h_0,\cdot]$ degree-preserving), have degree
$\le3n-2(n-1)=n+2$; Moyal corrections lower degree by even amounts,
fixing the parity $\equiv n\bmod2$.  A kernel monomial with transfer
$(d_1,d_2)$, $d_1d_2<0$, requires degree $\ge|d_1|+|d_2|$ with matching
per-mode parities, giving the candidate sets of
Theorem~\ref{thm:selection}.

\section{Field conventions, the source formula, and truncation
completeness}\label{app:field}

\emph{Conventions.}  Finite volume; discrete momenta $k$ (lattice units
with $m_L=4$ in Theorem~\ref{thm:sector}); modes $(b,k)$ with
$[a_{b,k},a^{\dg}_{b',k'}]=\delta_{bb'}\delta_{kk'}$; unit-norm
external states; $g=1$; deformation parameter $\zeta$ multiplying
$H_3$, $\zeta^2$ multiplying $H_4$.  Positive-frame legs from the
branch eigenvectors:
$u$-legs $(a_++a_+^{\dg})/\sqrt{2\w_+\delta}$ (healthy) and
$(a_--a_-^{\dg})/\sqrt{2\w_-\delta}$ (ghost, quarter-turn);
$v$-legs $\rho_\pm$ times $u$-legs, $\rho_\pm=\mp\delta/(2g)$.
Interactions $H_3=-g\!\int\!uv^2$, $H_4=-\tfrac g2\!\int\!u^2v^2$ are
pure potentials.  The transported fourth-order field of
Section~\ref{sec:firstfield} is fixed by matching the free two-point
function $[e^{-i\w_1t}/2\w_1-e^{-i\w_2t}/2\w_2]/\delta$ of
\cite{paperIII}; the even-ghost rule follows because a monomial with
$g$ ghost legs, $c_g$ of them creators, has $v$-coefficient
$\propto i^{\,\#\mathrm{legs}}(-1)^{c_g}$ and
$v^{\dg}$-coefficient $\propto$ the conjugate with
$c_g\to g-c_g$, whose difference vanishes for odd $g$.

\emph{Source formula.}  With $E_{\mathrm{in}}=E_{\mathrm{out}}=E$ and
$R_1$ solved off shell,
$\langle\mathrm{out}|[R_1,w]|\mathrm{in}\rangle
=\sum_n[(v^{\dg}-v)_{\mathrm{out},n}w_{n,\mathrm{in}}
+w_{\mathrm{out},n}(v^{\dg}-v)_{n,\mathrm{in}}]/(E-E_n)$ with
$w=v+v^{\dg}$; the algebraic identity
$(v^{\dg}-v)w+w(v^{\dg}-v)=2(v^{\dg}v^{\dg}-vv)$ reduces this to
\eqref{eq:sourceME}.

\emph{Truncation completeness (proof of Lemma~\ref{lem:trunc}).}
Second-order chains pass through 1-, 3-, or 5-particle intermediates.
1-particle: momentum $k_a+k_b$ fixed.  3-particle: the first vertex
annihilates one external and creates a pair $(q,\,k_x-q)$; for the
second vertex to reach the 2-particle out-state, the surviving created
particle must be an out-particle, forcing
$q\in\{k_c,k_d,k_x-k_c,k_x-k_d\}$; annihilating both created particles
instead returns the in-momenta, giving zero overlap with a distinct
out-state.  5-particle: the created triple $(p,q,r)$, $p+q+r=0$, must
contain both out-particles, fixing the third momentum.  Hence all
internal momenta lie in the $s,t,u$-generated set, and enlarging the
momentum cutoff cannot change the element --- as the exact agreement at
$K_{\max}=4,6$ confirms.

\section{The doubled confluent parity}\label{app:parity}

With change-of-basis matrix $C$ whose columns are $(c,d)$ in branch
coordinates, $P_\delta^{(c,d)}=C^{-1}\operatorname{diag}(1,-1)C
=\left(\begin{smallmatrix}0&2/\delta\\ \delta/2&0\end{smallmatrix}\right)$.
For the doubled space take $C_A=C$ and the oppositely oriented
$C_B$ (columns $c_B=(b_+-b_-)/2$, $d_B=(b_++b_-)/\delta$); the cross
map $\kappa(b_\pm^A)=\pm b_\pm^B$ has matrix
$T^{-1}\kappa_{\mathrm{branch}}T$ with $T=\operatorname{diag}(C_A,C_B)$,
which evaluates to the $\delta$-independent exchange
$(c_A,d_A)\leftrightarrow(c_B,d_B)$, squaring to the identity.

\section{Verification checks and their classification}
\label{app:checks}

\begin{center}
\small
\begin{tabular}{l l l}
\hline
Check & Supports & Type\\
\hline
ID1--ID5 & Theorem~\ref{thm:first} & exact symbolic\\
ID6--ID8 & Theorem~\ref{thm:secondgen} & exact symbolic\\
ID9a--b & Theorem~\ref{thm:threeone} & exact symbolic\\
ID10, O3e & Comp.\ Prop.~\ref{cprop:scaling} & high-precision numeric\\
SR1--SR3 & Theorem~\ref{thm:selection} & exact combinatorial\\
O3a--O3d & Theorem~\ref{thm:third} & exact symbolic\\
PT1--PT2 & Theorem~\ref{thm:complexpair} (structure) & exact symbolic\\
(new) Sturm count & Theorem~\ref{thm:complexpair} (pair) &
exact rational\\
PT3 & Theorem~\ref{thm:complexpair} (truncated op.) & numeric\\
PT4 & broken-PT tongues & numeric (effective matrix)\\
PS1--PS2 & Proposition~\ref{prop:kallen} & exact symbolic\\
PS-A--PS-G & Theorem~\ref{thm:firstfield}, Comp.\ Prop. &
exact symbolic + numeric\\
PS-H & Jordan-chain lemma & exact symbolic\\
TF1--TF7 & Theorem~\ref{thm:twofield}, Prop.~\ref{prop:sectors} &
exact symbolic\\
SO1--SO7 & Theorem~\ref{thm:sector} (numerics) & exact + numeric\\
HX1 & Theorem~\ref{thm:sector} (exact value) & exact symbolic\\
HX2--HX3 & Theorem~\ref{thm:confluent} & exact symbolic\\
(new) PT proposition & Proposition~\ref{prop:pt} & exact symbolic\\
(new) $\w_2^{-13/2}$ ratio & Theorem~\ref{thm:third} & exact symbolic\\
(new) $R_1$ Laurent & Theorem~\ref{thm:jordanexact} & exact symbolic\\
DQ1--DQ2, DQ9 & $O(1,1)$ polar form, current, Ward & exact symbolic\\
DQ3--DQ6 & doubled pairing, graph, pseudo-unitarity & numeric\\
DQ7--DQ8 & Comp.\ Prop.~\ref{cprop:krein} & exact symbolic\\
FO1--FO9 & Comp.\ Prop.~\ref{cprop:fiveone},
Conj.~\ref{conj:hierarchy} & exact symbolic\\
ON1--ON4 & Lemma~\ref{lem:chargenull},
Comp.\ Prop.~\ref{cprop:embedding} & exact symbolic + numeric\\
\hline
\end{tabular}
\end{center}

Scripts, in \texttt{symbolic/}, each printing \texttt{ALL PASS}:\\
\texttt{verify\_interaction\_deformation.py} \quad
\texttt{verify\_interaction\_order3.py}\\
\texttt{verify\_pt\_breaking.py} \quad
\texttt{verify\_perfect\_square.py} \quad
\texttt{verify\_two\_field.py}\\
\texttt{verify\_sector\_obstruction.py} \quad
\texttt{verify\_hardening.py}\\
\texttt{verify\_doubled\_theory.py} \quad
\texttt{verify\_51\_order4.py}\\
\texttt{verify\_obstruction\_null.py}

\end{document}
